\section{Limitations}\label{sec:limitations}

\begin{itemize}\setlength{\itemsep}{0pt}\setlength{\parskip}{0pt}
\item FX comes from one market maker, and its original evaluation period was inspected before conditional methods were specified.
\item The earlier-period comparison and full-prefix backcast are retrospective, with no previously unseen prospective financial evaluation.
\item Weekly reports aggregate positions across dealers; related categories and the CFTC survivor universe do not provide independent dealer replications. Net-position peaks summarize exposure; gross portfolios, hedging arrangements and economic capital requirements remain unidentified.
\item Downloaded economic-date vintages omit historical first-release information, reporting delays and unrecorded within-interval excursions.
\item The comparisons identify neither causal effects nor individual-feature effects, and fitted limits receive no coverage guarantee under this dependence.
\item Costs illustrate fixed preferences; trading returns, funding gains and institutional intervention costs are not measured.
\item The finite-penalty hindsight bound is diagnostic only; validation-length and quantile-sensitivity checks are post-result analyses with numerical failures retained.
\item Only 6 of 10 high-quantile compatibility fits pass, so the rejected repair leaves both FX comparisons unavailable.
\end{itemize}
Supplement Section S36 retains the original technical qualifications on simulations, aggregation and forecast comparison. A consolidated register in Supplement Section S40 retains supported, unsupported and untestable prewritten screens.

\FloatBarrier

\section{Conclusion}\label{sec:conclusion}

Historical rules are hard to beat. For RQ1, matched FX regression gains recur across periods, but a shorter rolling reference changes the ranking (Table~\ref{tab:fx}). The learning curve also gives back part of its gain at the longest history, linking conditional performance to available history rather than universal model advantage (Figure~\ref{fig:v11-history}).

For RQ2, longer validation reduces full regression medians in both weekly reporting systems, but across-arm differences need not describe typical paired improvements (Table~\ref{tab:v13-validation}). These dealer aggregates remain related tasks, and the horizon-six New York Fed extension also moves the validation endpoint. The benchmark-default protocol cuts weekly median excess to 0.00\% and preserves FX gains, yet 8 of 60 cases still lose by over ten percent, with a worst excess of 80.98\% (Section~\ref{sec:v13-rp}). This failed tail criterion limits the protocol's risk-management recommendation, and the rejected high-quantile repair leaves its original judgments unresolved.

For RQ3, higher reserve charges and lower fixed breach charges generally favor expanded regression in CFTC, with local exceptions (Figure~\ref{fig:v10-cost}). Lower pinball scores can therefore coexist with more breaches and different cost rankings, requiring joint reporting of reserve amounts, breach counts and exceedance amounts. The costs are stipulated preferences rather than observed institutional costs, and the financial evidence remains retrospective.

\FloatBarrier

